\documentclass[10pt,a4paper]{amsart}
\usepackage[margin=19mm]{geometry}
\usepackage{amsmath,amssymb,mathtools}
\usepackage[scr=rsfs,cal=euler]{mathalfa}
\usepackage{xfrac}
\usepackage[unicode,hidelinks,bookmarks=false]{hyperref}
\newcommand{\ie}{i.e.\ }
\newcommand{\cf}{cf.\ }
\newcommand{\resp}{respectively\ }
\newcommand{\Subsection}{\subsection}
\providecommand{\nobreakdash}{\nobreak\hbox{-}}
\setlength{\emergencystretch}{2em}
\multlinegap0pt
\usepackage{textcomp}
\usepackage{mathtools,mathrsfs}
\usepackage{enumitem}
\usepackage{booktabs}
\usepackage{array}
\hypersetup{
  colorlinks=true,
  linkcolor=blue,
  citecolor=blue,
  urlcolor=blue
}
\newtheorem{theorem}{Theorem}[section]
\newtheorem{proposition}[theorem]{Proposition}
\newtheorem{definition}[theorem]{Definition}
\newtheorem{lemma}[theorem]{Lemma}
\newtheorem{remark}[theorem]{Remark}
\newtheorem{assumption}[theorem]{Assumption}
\newtheorem{corollary}[theorem]{Corollary}
\newtheorem{example}[theorem]{Example}
\newcommand{\R}{\mathbb{R}}
\newcommand{\Om}{\Omega}
\newcommand{\pa}{\partial}
\newcommand{\cl}{\overline}
\newcommand{\grad}{\nabla}
\newcommand{\diver}{\operatorname{div}}
\newcommand{\tr}{\operatorname{tr}}
\newcommand{\id}{\mathrm{id}}
\newcommand{\dd}{\,\mathrm{d}}
\newcommand{\F}{\mathcal{F}}
\newcommand{\X}{\mathcal{X}}
\newcommand{\Diss}{\mathcal{D}}
\newcommand{\Bnd}{\mathcal{B}}
\newcommand{\Src}{\mathcal{S}}
\newcommand{\ThermCont}{\mathbf{ThermCont}}
\newcommand{\ThermAdm}{\mathbf{ThermCont}^{\mathrm{adm}}}
\newcommand{\ThermLin}{\mathbf{ThermLin}}
\newcommand{\ThermDisc}{\mathbf{ThermDisc}}
\newcommand{\lin}{\mathrm{lin}}
\newcommand{\pair}[2]{\left\langle #1,#2\right\rangle}
\begin{document}
\title[Thermodynamic Structure and Composition in Nonlinear Convect]{Thermodynamic Structure and Composition in Nonlinear Convection-Diffusion}
\author{J. J. Segura}
\date{}
\maketitle
\noindent\emph{Source and licence.} Thermodynamic Structure and Composition in Nonlinear Convection-Diffusion. Version arXiv:2607.07829v1 (CC BY 4.0). This material is distributed under the Creative Commons Attribution 4.0 International licence, http://creativecommons.org/licenses/by/4.0/, as recorded in the arXiv OAI licence metadata for this exact version and shown on the abstract page https://arxiv.org/abs/2607.07829v1. Full rights notice: licenses/arxiv-2607.07829-CC-BY-4.0.txt.\par\smallskip
\noindent\textit{Editorial scope.} Counted unit: the original \texttt{theorem} environment \texttt{thm:interconnection} at source lines 814--835 together with its complete proof at lines 837--847; the delivered document keeps source lines 539--848 so that every referenced label and every citation resolves inside it. Native editable source is the paper's own concatenated TeX; the AMS wrapper is editorial formatting only. No publisher class, upstream script or unrelated artwork is redistributed.
\section{Restriction, local-to-global reconstruction, and interconnection}

This section sharpens the compositional core of the manuscript. The key point is that
thermodynamic consistency is not only stable under passing to a subdomain or under
coupling two open subsystems; it also admits a genuine local-to-global
interpretation. Local balances on a compatible family of subdomains reconstruct the
global balance exactly, because the internal interface exchanges cancel pairwise.
This is the first place where the paper's categorical thesis acquires a distinctly
geometric form.

\subsection{Admissible subdomains and local thermodynamic data}

Let \(\mathsf{Dom}(\Om)\) denote the category whose objects are Lipschitz subdomains
\(U\subseteq \Om\), and whose morphisms are inclusions \(V\hookrightarrow U\). The
category is a poset category under set inclusion.

For each \(U\in \mathsf{Dom}(\Om)\), write
\[
\Gamma_U^{\mathrm{ext}}:=\pa U\cap \pa\Om,
\qquad
\Gamma_U^{\mathrm{int}}:=\pa U\cap \Om.
\]
If \(u\) is a smooth admissible trajectory on \(\Om\), define the restricted local
thermodynamic data on \(U\) by
\begin{align}
\F_U[u(t)]
&:=
\int_U \psi(x,u(x,t))\,\dd x,\label{eq:F-U}\\
\Diss_U[u(t)]
&:=
\int_U \sigma[u(t)]\,\dd x,\label{eq:D-U}\\
\Src_U[u(t)]
&:=
\int_U \mu[u(t)]\cdot R[u(t)]\,\dd x,\label{eq:S-U}\\
\Bnd_U^{\mathrm{ext}}[u(t)]
&:=
-\int_{\Gamma_U^{\mathrm{ext}}} q[u(t)]\cdot n_U\,\dd S,\label{eq:B-U-ext}\\
\Bnd_U^{\mathrm{int}}[u(t)]
&:=
-\int_{\Gamma_U^{\mathrm{int}}} q[u(t)]\cdot n_U\,\dd S,\label{eq:B-U-int}
\end{align}
where \(q=q^{\mathrm{rev}}+q^{\mathrm{diss}}\) and \(n_U\) is the outward unit normal
on \(\pa U\). We also set
\begin{equation}\label{eq:B-U-total}
\Bnd_U[u]:=\Bnd_U^{\mathrm{ext}}[u]+\Bnd_U^{\mathrm{int}}[u].
\end{equation}

\begin{definition}[Category of local thermodynamic data]
Let \(\mathsf{Data}\) denote the category whose objects are tuples
\[
(\F,\Diss,\Src,\Bnd^{\mathrm{ext}},\Bnd^{\mathrm{int}})
\]
associated with a Lipschitz subdomain, and whose morphisms are the restriction maps
induced by admissible inclusions of subdomains.
\end{definition}

\begin{definition}[Local thermodynamic assignment]
The \emph{local thermodynamic assignment} associated with a global system
\(\mathsf{T}\in\ThermAdm\) is the functor
\[
\mathbb{T}_{\mathsf{T}}:\mathsf{Dom}(\Om)^{\mathrm{op}}\to \mathsf{Data},
\qquad
U\mapsto
\bigl(\F_U,\Diss_U,\Src_U,\Bnd_U^{\mathrm{ext}},\Bnd_U^{\mathrm{int}}\bigr).
\]
\end{definition}

\begin{proposition}[Contravariant locality structure]\label{prop:local-assignment}
The assignment \(\mathbb{T}_{\mathsf{T}}\) is contravariant with respect to
admissible inclusions. In particular, if \(V\subseteq U\subseteq \Om\), then the
thermodynamic data on \(V\) are obtained from those on \(U\) by restriction of the
underlying fields and traces.
\end{proposition}

\begin{proof}
All local quantities are defined by integration or trace over the chosen subdomain.
If \(V\subseteq U\), then the restricted field \(u|_V\) determines \(\F_V\),
\(\Diss_V\), \(\Src_V\), and the corresponding external and internal boundary terms
by the same formulas. Compatibility with identity maps and composition of inclusions
is immediate.
\end{proof}

\begin{remark}
Proposition~\ref{prop:local-assignment} is only a first structural layer. It says
that thermodynamic data localize contravariantly over subdomains, but it does not
yet impose a sheaf axiom. The gluing content appears only once interface
compatibility is taken into account.
\end{remark}

\subsection{Restriction to a single subdomain}

\begin{theorem}[Restriction theorem]\label{thm:restriction}
Let \(\mathsf{T}\in\ThermAdm\), let \(u\) be a smooth admissible trajectory of
\(\mathsf{T}\), and let \(U\in\mathsf{Dom}(\Om)\). Then
\begin{equation}\label{eq:balance-U}
\frac{\dd}{\dd t}\F_U[u(t)]
=
-\Diss_U[u(t)]
+\Bnd_U^{\mathrm{ext}}[u(t)]
+\Bnd_U^{\mathrm{int}}[u(t)]
+\Src_U[u(t)].
\end{equation}
\end{theorem}

\begin{proof}
Integrate the local thermodynamic balance
\[
\partial_t \psi(\cdot,u) + \diver q[u]
=
-\sigma[u] + \mu[u]\cdot R[u]
\]
over \(U\). The divergence term splits over \(\Gamma_U^{\mathrm{ext}}\) and
\(\Gamma_U^{\mathrm{int}}\), giving exactly the terms
\eqref{eq:B-U-ext}--\eqref{eq:B-U-int}. Rearrangement yields
\eqref{eq:balance-U}.
\end{proof}

\begin{remark}
Theorem~\ref{thm:restriction} shows that restriction does not break the
thermodynamic law. It merely separates boundary exchange into two visible pieces:
external exchange with the environment and internal exchange across the newly exposed
interface.
\end{remark}

\subsection{Compatible partitions and local-to-global reconstruction}

\begin{definition}[Compatible finite partition]\label{def:compatible-partition}
A finite family \(\{\Om_i\}_{i=1}^N\subseteq \mathsf{Dom}(\Om)\) is called a
\emph{compatible finite partition} of \(\Om\) if:
\begin{enumerate}[label=(\roman*),leftmargin=1.1cm]
    \item \(\Om=\bigcup_{i=1}^N \cl{\Om_i}\);
    \item the interiors of the \(\Om_i\) are pairwise disjoint;
    \item each \(\Om_i\) is Lipschitz;
    \item for each \(i\neq j\), the interface
    \[
    \Gamma_{ij}:=\pa\Om_i\cap \pa\Om_j\cap \Om
    \]
    is either empty or a Lipschitz hypersurface on which the outer normals satisfy
    \[
    n_j=-n_i
    \qquad \text{a.e. on }\Gamma_{ij}.
    \]
\end{enumerate}
\end{definition}

For such a partition, define the local data on each \(\Om_i\) by
\[
\F_i:=\F_{\Om_i},\qquad
\Diss_i:=\Diss_{\Om_i},\qquad
\Src_i:=\Src_{\Om_i},
\]
and write
\begin{align}
\Bnd_i^{\mathrm{ext}}[u]
&=
-\int_{\pa\Om_i\cap \pa\Om} q[u]\cdot n_i\,\dd S,\label{eq:Biext}\\
\Bnd_i^{\mathrm{int}}[u]
&=
-\sum_{j\neq i}\int_{\Gamma_{ij}} q[u]\cdot n_i\,\dd S.\label{eq:Biint}
\end{align}

\begin{lemma}[Pairwise interface cancellation]\label{lem:pairwise-cancel}
Let \(\{\Om_i\}_{i=1}^N\) be a compatible finite partition of \(\Om\), and let \(u\)
be a smooth admissible trajectory. Then, for each nonempty interface \(\Gamma_{ij}\),
\begin{equation}\label{eq:pairwise-cancel}
-\int_{\Gamma_{ij}} q[u]\cdot n_i\,\dd S
-
\int_{\Gamma_{ij}} q[u]\cdot n_j\,\dd S
=
0.
\end{equation}
Consequently,
\begin{equation}\label{eq:sum-Bint-zero}
\sum_{i=1}^N \Bnd_i^{\mathrm{int}}[u]=0.
\end{equation}
\end{lemma}

\begin{proof}
On \(\Gamma_{ij}\), the compatible partition assumption gives \(n_j=-n_i\). Hence
\[
q[u]\cdot n_j = -\,q[u]\cdot n_i
\qquad \text{a.e. on }\Gamma_{ij},
\]
which proves \eqref{eq:pairwise-cancel}. Summing over all interfaces yields
\eqref{eq:sum-Bint-zero}.
\end{proof}

\begin{theorem}[Local-to-global reconstruction]\label{thm:local-to-global}
Let \(\mathsf{T}\in\ThermAdm\), let \(u\) be a smooth admissible trajectory, and let
\(\{\Om_i\}_{i=1}^N\) be a compatible finite partition of \(\Om\). Then:
\begin{align}
\F[u] &= \sum_{i=1}^N \F_i[u],\label{eq:F-sum}\\
\Diss[u] &= \sum_{i=1}^N \Diss_i[u],\label{eq:D-sum}\\
\Src[u] &= \sum_{i=1}^N \Src_i[u],\label{eq:S-sum}\\
\Bnd[u] &= \sum_{i=1}^N \Bnd_i^{\mathrm{ext}}[u],\label{eq:B-sum-ext}
\end{align}
and summing the local balances \eqref{eq:balance-U} over \(i\) recovers the global
balance
\begin{equation}\label{eq:global-recovered}
\frac{\dd}{\dd t}\F[u(t)]
=
-\Diss[u(t)] + \Bnd[u(t)] + \Src[u(t)].
\end{equation}
\end{theorem}

\begin{proof}
The identities \eqref{eq:F-sum}--\eqref{eq:S-sum} follow from additivity of
integrals over a partition. Summing \eqref{eq:balance-U} over \(i\) gives
\[
\frac{\dd}{\dd t}\sum_{i=1}^N \F_i[u(t)]
=
-\sum_{i=1}^N \Diss_i[u(t)]
+\sum_{i=1}^N \Bnd_i^{\mathrm{ext}}[u(t)]
+\sum_{i=1}^N \Bnd_i^{\mathrm{int}}[u(t)]
+\sum_{i=1}^N \Src_i[u(t)].
\]
Use \eqref{eq:F-sum}--\eqref{eq:S-sum}, the interface cancellation
\eqref{eq:sum-Bint-zero}, and the fact that the sum of the external pieces is the
global boundary exchange \eqref{eq:B-sum-ext}. This yields
\eqref{eq:global-recovered}.
\end{proof}

\begin{remark}[Local-to-global principle]
Theorem~\ref{thm:local-to-global} is the clearest expression of the paper's central
claim. The second law is not merely stable under restriction; it is reconstructible
from compatible local balances. Thermodynamic consistency therefore behaves as a
local-to-global invariant.
\end{remark}

\subsection{Open subsystems and power-conserving gluing}

\begin{definition}[Open local subsystem]
An \emph{open local subsystem} on a Lipschitz domain \(U\subseteq \Om\) is an object
\[
\mathsf{T}_U
=
(U,\X_U,\F_U,\Diss_U,\Src_U,
Y_{e,U}^{\Gamma},Y_{f,U}^{\Gamma},e_U^\Gamma,f_U^\Gamma)
\]
equipped with a local balance
\begin{equation}\label{eq:local-open-balance}
\frac{\dd}{\dd t}\F_U[u_U(t)]
=
-\Diss_U[u_U(t)]
+\Bnd_U^{\mathrm{ext}}[u_U(t)]
+\Bnd_U^{\Gamma}[u_U(t)]
+\Src_U[u_U(t)],
\end{equation}
where
\[
\Bnd_U^{\Gamma}[u_U]:=
\pair{e_U^\Gamma[u_U]}{f_U^\Gamma[u_U]}_{Y_U^\Gamma}.
\]
\end{definition}

\begin{definition}[Interface interconnection relation]
Let \(\mathsf{T}_{\Om_i}\) and \(\mathsf{T}_{\Om_j}\) be open local subsystems sharing an
interface \(\Gamma_{ij}\). An \emph{interface interconnection relation} is a subset
\[
\mathfrak{I}_{ij}
\subset
Y_{e,\Om_i}^{\Gamma_{ij}}\times Y_{f,\Om_i}^{\Gamma_{ij}}
\times
Y_{e,\Om_j}^{\Gamma_{ij}}\times Y_{f,\Om_j}^{\Gamma_{ij}}
\]
such that every quadruple
\((\eta_i,\zeta_i,\eta_j,\zeta_j)\in \mathfrak{I}_{ij}\) satisfies
\begin{equation}\label{eq:power-ij}
\pair{\eta_i}{\zeta_i}_{Y_{\Om_i}^{\Gamma_{ij}}}
+
\pair{\eta_j}{\zeta_j}_{Y_{\Om_j}^{\Gamma_{ij}}}
=0.
\end{equation}
\end{definition}


\begin{theorem}[Interconnection theorem]\label{thm:interconnection}
Let \(\{\Om_i\}_{i=1}^N\) be a compatible finite partition of \(\Om\), and let each
\(\Om_i\) carry an open local subsystem \(\mathsf{T}_{\Om_i}\) satisfying
\eqref{eq:local-open-balance}. Assume that on every nonempty interface \(\Gamma_{ij}\)
the subsystem trajectories satisfy a power-conserving interconnection relation
\(\mathfrak{I}_{ij}\). Then the composite system with total storage, dissipation,
source, and external boundary exchange
\begin{align}
\F_\#[u_1,\dots,u_N] &:= \sum_{i=1}^N \F_i[u_i],\label{eq:Fsharp-local}\\
\Diss_\#[u_1,\dots,u_N] &:= \sum_{i=1}^N \Diss_i[u_i],\label{eq:Dsharp-local}\\
\Src_\#[u_1,\dots,u_N] &:= \sum_{i=1}^N \Src_i[u_i],\label{eq:Ssharp-local}\\
\Bnd_\#[u_1,\dots,u_N] &:= \sum_{i=1}^N \Bnd_i^{\mathrm{ext}}[u_i]\label{eq:Bsharp-local}
\end{align}
satisfies
\begin{equation}\label{eq:sharp-local-balance}
\frac{\dd}{\dd t}\F_\#[u_1(t),\dots,u_N(t)]
=
-\Diss_\#[u_1(t),\dots,u_N(t)]
+\Bnd_\#[u_1(t),\dots,u_N(t)]
+\Src_\#[u_1(t),\dots,u_N(t)].
\end{equation}
\end{theorem}

\begin{proof}
Sum the local open balances \eqref{eq:local-open-balance} over \(i\). The total
interface contribution is
\[
\sum_{i=1}^N \Bnd_i^{\Gamma}[u_i(t)].
\]
Group these terms by interfaces \(\Gamma_{ij}\). By the power-conserving relation
\eqref{eq:power-ij}, each pair of interface contributions cancels. The remaining
terms are exactly the sums \eqref{eq:Fsharp-local}--\eqref{eq:Bsharp-local}, which
yields \eqref{eq:sharp-local-balance}.
\end{proof}


\end{document}
